\chapter{Base, roll et livraison}\label{ch:m1:basis-roll-and-delivery}

L'indice est à 6\,000 et son contrat à terme de décembre se négocie à 6\,043. Un
commentateur lit dans ces quarante-trois points un signe d'optimisme. Ils
sont des intérêts~: l'acheteur du contrat à terme garde 300\,000~\$ à la banque pendant trois mois au lieu de
payer des actions, il en tire environ 3\,185~\$ et renonce à environ 1\,050~\$ de
dividendes, et la différence, divisée par le multiplicateur, vaut quarante-trois
points à la décimale près. La prime aura fondu à zéro d'ici le troisième vendredi de
décembre, quoi que fasse le marché. Le 20 avril 2020, le même mécanisme a joué
à l'envers sur le pétrole brut~: les détenteurs d'un contrat à un jour de l'échéance, face
à la livraison de barils qu'ils ne pouvaient pas stocker, ont payé à d'autres 37.63~\$ le baril pour se défaire des
contrats. Un contrat à terme est lié à son sous-jacent par un arbitrage et par
une procédure de livraison, et les deux ont un coût. Ce chapitre chiffre ce lien.

\section{Coût de portage et fair value}

\begin{definition}[Base]\label{def:m1:basis-roll-and-delivery:basis}
La \emph{base}\index{base} d'un contrat à terme est la différence entre son prix
et le prix au comptant de son sous-jacent, ici le contrat à terme moins le comptant. Elle converge vers
zéro à l'échéance, car à l'échéance le contrat à terme \emph{est} le sous-jacent.
\end{definition}

\begin{definition}[Coût de portage et fair value]\label{def:m1:basis-roll-and-delivery:carry}
Le \emph{coût de portage}\index{coût de portage} d'un actif est ce qu'il coûte de
le détenir jusqu'à une date ultérieure~: le financement, auquel s'ajoutent le stockage et l'assurance pour une
matière première, diminué des revenus qu'il verse. La \emph{fair value}\index{fair value}
d'un contrat à terme est le prix au comptant augmenté du coût de portage jusqu'à l'échéance~: le cours du
contrat à terme auquel l'achat de l'actif et la vente du contrat rapportent exactement le
taux de financement.
\end{definition}

\begin{proposition}[Fair value d'un contrat à terme sur indice]\label{prop:m1:basis-roll-and-delivery:fv}
Soit $S$ l'indice, $r$ le taux de financement (simple, exact/360), $T$ la
durée jusqu'à l'échéance et $d_i$ les dividendes, en points d'indice, versés aux dates $t_i
\le T$. Alors
\[
F^* \;=\; S\,(1 + rT) \;-\; \sum_i d_i\,\bigl(1 + r\,(T - t_i)\bigr).
\]
\end{proposition}
\begin{proof}
Emprunter $S$, acheter le panier, vendre un contrat à terme à $F$. Placer chaque dividende jusqu'à
l'échéance. À l'échéance, livrer la valeur du panier contre le contrat à terme~: la trésorerie est
$F + \sum_i d_i(1 + r(T-t_i)) - S(1+rT)$ quoi que fasse l'indice. Sans
mise de fonds initiale ni risque, ce montant doit être nul.
\end{proof}

\begin{example}[Quarante-trois points]\label{ex:m1:basis-roll-and-delivery:43}
Le 18 septembre 2026, avec $S = 6\,000$, $r = 4.2\%$, échéance le 18 décembre (91
jours) et des dividendes de 6.0, 9.5 et 5.5 points versés les 15 octobre, 16
novembre et 10 décembre~: les intérêts valent $6\,000 \times 0.042 \times 91/360 =
63.70$ points, les dividendes valent 21.09 à l'échéance, et $F^* = 6\,042.61$.
Un cours de marché de 6\,047.20 correspond à un \emph{taux de financement
implicite}\index{taux de financement implicite} de 4.50~\%~: le \omterm{def:m1:delta-one-instruments:swap}{spread de financement}
du \cref{ch:m1:delta-one-instruments}, 30 points de base, vu de l'autre côté.
\end{example}

\begin{omfigure}
\begin{tikzpicture}
\begin{axis}[width=11cm,height=4.8cm,
  xlabel={jours avant l'échéance},ylabel={base juste (points d'indice)},x dir=reverse,
  tick label style={font=\small},label style={font=\small},
  xmin=0,xmax=91,ymin=0,ymax=48,grid=major,grid style={gray!25},table/col sep=comma]
\addplot[omDef,very thick,const plot] table[x=days_to_expiry,y=basis]{figdata/markets-1/21-basis-roll-and-delivery/basis.csv};
\node[font=\footnotesize,anchor=west] at (axis cs:63,33) {dividende 6.0};
\node[font=\footnotesize,anchor=west] at (axis cs:31,27) {9.5};
\node[font=\footnotesize,anchor=west] at (axis cs:7.5,12) {5.5};
\end{axis}
\end{tikzpicture}

\omcaption{La \omterm{def:m1:basis-roll-and-delivery:basis}{base} juste de l'\cref{ex:m1:basis-roll-and-delivery:43} avec
l'indice maintenu à 6\,000~: les intérêts s'accumulent à raison de 0.7 point par jour, et la \omterm{def:m1:basis-roll-and-delivery:basis}{base}
bondit \emph{vers le haut} à chaque versement de dividende, puisque ce dividende
n'est plus à venir. Une série de \omterm{def:m1:basis-roll-and-delivery:basis}{bases} présente ce profil en dents de scie avant toute
transaction. Données~: la construction du chapitre.}\label{fig:m1:basis-roll-and-delivery:basis}
\end{omfigure}

\section{Arbitrage sur indice}

\begin{proposition}[La bande d'arbitrage]\label{prop:m1:basis-roll-and-delivery:band}
Avec des coûts de transaction à sens unique $c_S$ sur le panier et $c_F$ sur le contrat à terme (en
fractions de la valeur), un \omterm{def:m1:delta-one-instruments:swap}{spread de financement} $\phi$ pour l'arbitragiste lorsqu'il est
acheteur du panier et des \omterm{def:m1:financing:seclending}{frais d'emprunt} $\beta$ lorsqu'il est vendeur à découvert, le contrat à terme peut se négocier
sans arbitrage en aucun point de
\[
\Bigl[F^* - S\bigl(2(c_S + c_F) + \beta T\bigr),\;\; F^* + S\bigl(2(c_S + c_F) + \phi T\bigr)\Bigr].
\]
La bande est plus étroite pour les acteurs à coûts plus faibles, se resserre à mesure que l'échéance
approche, et elle est asymétrique lorsque emprunter le panier coûte plus cher que
le financer.
\end{proposition}

\begin{example}[Une bande de dix-sept points]\label{ex:m1:basis-roll-and-delivery:band}
Avec $c_S = 3$ et $c_F = 0.5$ point de base, $\phi = 15$ et $\beta = 40$ points de
base par an, celle de l'\cref{ex:m1:basis-roll-and-delivery:43} s'étend de
$F^* - 10.27$ à $F^* + 6.47$~: soit de 6\,032.35 à 6\,049.09. À 6\,050.50, le
contrat à terme est surévalué de 7.89 points et la transaction, acheter le panier et vendre le contrat à terme,
verrouille 1.41 points, soit 70~\$ par contrat, conservée jusqu'à l'échéance. La plupart des arbitrages sur indice
ne sont pas conservés jusqu'à l'échéance~: la position est dénouée lorsque l'écart de valorisation s'inverse,
ce qui rapporte l'aller-retour de l'écart de valorisation et supporte les coûts de transaction deux
fois de plus.
\end{example}

\begin{omfigure}
\begin{tikzpicture}
\begin{axis}[width=11cm,height=4.8cm,
  xlabel={temps (pas arbitraires)},ylabel={future moins fair value},
  tick label style={font=\small},label style={font=\small},
  xmin=0,xmax=400,ymin=-13,ymax=9,grid=major,grid style={gray!25},table/col sep=comma]
\addplot[omDef,thick] table[x=t,y=mispricing]{figdata/markets-1/21-basis-roll-and-delivery/mispricing.csv};
\addplot[omThm,thick,dashed] table[x=t,y=upper]{figdata/markets-1/21-basis-roll-and-delivery/mispricing.csv};
\addplot[omThm,thick,dashed] table[x=t,y=lower]{figdata/markets-1/21-basis-roll-and-delivery/mispricing.csv};
\end{axis}
\end{tikzpicture}

\omcaption{Un écart de prix simulé à l'intérieur de la bande de
l'\cref{ex:m1:basis-roll-and-delivery:band} (en pointillés, en points d'indice). La borne
supérieure est plus proche que la borne inférieure, car financer un panier acheteur coûte
moins cher que d'emprunter un panier vendeur, et elle est atteinte plus souvent. Données~:
simulation du tutoriel.}\label{fig:m1:basis-roll-and-delivery:mispricing}
\end{omfigure}

\section{Le roll}

\begin{definition}[Roll]\label{def:m1:basis-roll-and-delivery:roll}
Effectuer le \emph{roll}\index{roll} d'une position à terme consiste à la clôturer sur le contrat arrivant à échéance
et à la rouvrir sur un contrat plus éloigné, normalement sous la forme d'une seule transaction de \omterm{def:m1:matching-algorithms-and-implied-spreads:calendar}{calendar spread}
(\cref{def:m1:matching-algorithms-and-implied-spreads:calendar}). Le prix du
roll est le \omterm{def:m1:matching-algorithms-and-implied-spreads:calendar}{calendar spread}~; sa \emph{richesse} (richness) est l'excédent du
\omterm{def:m1:matching-algorithms-and-implied-spreads:calendar}{calendar spread} sur la différence des deux \omterm{def:m1:basis-roll-and-delivery:carry}{fair values}, généralement exprimé en
taux annualisé.
\end{definition}

Pour un indice boursier, le \omterm{def:m1:basis-roll-and-delivery:roll}{roll} tarifie trois mois de financement moins trois
mois de dividendes. Un porteur acheteur qui effectue son \omterm{def:m1:basis-roll-and-delivery:roll}{roll} à un spread supérieur de 4.55 points à la
\omterm{def:m1:basis-roll-and-delivery:carry}{fair value} sur un indice de 6\,000 paie 30 points de base par an de plus que
l'indice de référence pour son effet de levier~; le porteur vendeur situé en face les encaisse.
Le \omterm{def:m1:basis-roll-and-delivery:roll}{roll} s'étale sur une semaine environ (\cref{fig:m1:futures-contracts-and-exchanges:roll}),
et comme tout le monde doit le faire, sa richesse constitue la mesure publique la plus nette
du prix du bilan.

\begin{definition}[Contango et backwardation]\label{def:m1:basis-roll-and-delivery:contango}
Une courbe de contrats à terme est en \emph{contango}\index{contango} lorsque les échéances plus lointaines
se traitent au-dessus des plus proches, et en \emph{backwardation}\index{backwardation}
lorsqu'elles se traitent en dessous.
\end{definition}

Pour un actif financier, la pente de la courbe est purement arithmétique~: taux moins rendement. Pour
une matière première, elle intègre le stockage et un \emph{rendement de convenance} (convenience yield), c'est-à-dire la valeur de
disposer du bien physique sous la main, valeur qui est élevée lorsque les stocks sont
faibles. La pente est un rendement~: une position longue reconduite en descendant une courbe en \omterm{def:m1:basis-roll-and-delivery:contango}{backwardation}
gagne cette pente même si le prix au comptant ne bouge jamais, et une position reconduite en montant une
courbe en \omterm{def:m1:basis-roll-and-delivery:contango}{contango} la paie.

\begin{omfigure}
\begin{tikzpicture}
\begin{axis}[width=11cm,height=4.8cm,
  xlabel={mois},ylabel={valeur de 1 investi},
  tick label style={font=\small},label style={font=\small},
  xmin=0,xmax=24,ymin=0.6,ymax=1.4,grid=major,grid style={gray!25},
  legend columns=2,legend style={font=\footnotesize,at={(0.5,-0.32)},anchor=north,draw=none,fill=none,/tikz/every even column/.append style={column sep=8pt}},
  table/col sep=comma]
\addplot[omThm,very thick] table[x=month,y=contango]{figdata/markets-1/21-basis-roll-and-delivery/rolled.csv};
\addplot[omDef,very thick,dashed] table[x=month,y=backwardation]{figdata/markets-1/21-basis-roll-and-delivery/rolled.csv};
\legend{deuxième mois 1.2~\% au-dessus du premier,deuxième mois 1.0~\% en dessous du premier}
\end{axis}
\end{tikzpicture}

\omcaption{Une position longue reconduite chaque mois pendant deux ans, alors que le prix au comptant
ne bouge jamais de 70. En \omterm{def:m1:basis-roll-and-delivery:contango}{contango}, elle perd 25~\%~; en \omterm{def:m1:basis-roll-and-delivery:contango}{backwardation}, elle gagne
27~\%. Un investisseur qui «~achète du pétrole~» par contrats à terme achète aussi la courbe.
Données~: simulation du tutoriel.}\label{fig:m1:basis-roll-and-delivery:rolled}
\end{omfigure}

\section{Livraison et règlement-livraison final}

\begin{definition}[Règlement en espèces et livraison physique]\label{def:m1:basis-roll-and-delivery:delivery}
À l'échéance, un contrat à terme \emph{réglé en espèces}\index{réglé en espèces} (cash-settled) verse une dernière
\omterm{def:m1:clearing-and-settlement:margin}{marge de variation} sur la \omterm{def:m1:basis-roll-and-delivery:basis}{base} d'un cours de règlement-livraison final, puis disparaît.
Un contrat à terme \emph{à livraison physique}\index{alivraisonphysique@à livraison physique} (physically delivered) oblige le
vendeur à livrer le sous-jacent, et l'acheteur à le recevoir et à le payer, selon les
règles de livraison de la bourse~: qualité, lieu, dates et, pour les contrats à terme sur obligations, le
choix parmi les titres livrables.
\end{definition}

\begin{definition}[Cours spécial d'ouverture]\label{def:m1:basis-roll-and-delivery:soq}
Un \emph{cours spécial d'ouverture}\index{cours spécial d'ouverture} (special opening quotation) est une
valeur de règlement-livraison finale d'un indice, calculée à partir des cours d'ouverture de chacun
de ses composants le jour du règlement-livraison. Ce n'est pas un niveau d'indice que
quiconque a vu~: les composants ouvrent à des moments différents, et le cours peut
se situer en dehors de la fourchette de la journée de l'indice publié.
\end{definition}

\begin{dated}{2026-09}{Comment trois contrats prennent fin}\label{dat:m1:basis-roll-and-delivery:ends}
\textbf{E-mini S\&P 500}~: la négociation du contrat arrivant à échéance s'arrête à
l'ouverture prévue de la New York Stock Exchange le troisième vendredi~; le règlement-livraison final
est un \omterm{def:m1:basis-roll-and-delivery:soq}{cours spécial d'ouverture} établi à partir des cours d'ouverture de ce
matin-là. \textbf{Contrat à terme sur l'EURO STOXX 50}~: la négociation s'arrête à midi le
troisième vendredi, sur la moyenne de l'indice au cours des dix dernières minutes.
\textbf{WTI brut}~: livraison physique à Cushing~; la négociation prend fin quelques jours
avant le début du mois de livraison, et un porteur encore acheteur à ce moment-là possède
du pétrole dans un hub de pipelines de l'Oklahoma.
\end{dated}

\begin{omfigure}
\begin{tikzpicture}[font=\small]
\draw[thick,-{Stealth[length=2.5mm]}] (0,0) -- (13.2,0);
\foreach \x/\d/\t in {1.0/jeudi 16:00/{dernière clôture des actions},5.2/vendredi 09:30/{fin de cotation des futures},11.0/{vendredi, plus tard}/{cours publié}}{
  \draw[thick] (\x,0.12) -- (\x,-0.12);
  \node[above=3pt,font=\footnotesize\bfseries] at (\x,0.12) {\d};
  \node[below=3pt,font=\footnotesize,align=center] at (\x,-0.12) {\t};}
\foreach \x in {5.5,5.7,5.8,6.0,6.1,6.4,6.6,7.1,7.9,9.2}{\draw[omThm,thick] (\x,0.22) -- (\x,0);}
\draw[decorate,decoration={brace,amplitude=5pt},omThm,thick] (5.45,0.95) -- (9.25,0.95) node[midway,above=6pt,font=\footnotesize,text=omThm,align=center] {les cours d'ouverture de chaque action~:\\le cours spécial est assemblé à partir de ceux-ci};
\draw[decorate,decoration={brace,amplitude=5pt,mirror},omDef,thick] (1.0,-0.95) -- (5.2,-0.95) node[midway,below=6pt,font=\footnotesize,text=omDef,align=center] {pendant la nuit~: les contrats à terme se négocient,\\pas les actions};
\end{tikzpicture}

\omcaption{L'échéance d'un contrat à terme sur indice \omterm{def:m1:basis-roll-and-delivery:delivery}{réglé en espèces} qui utilise un \omterm{def:m1:basis-roll-and-delivery:soq}{cours spécial
d'ouverture} (heure de New York). Entre le moment où les contrats à terme cessent
d'être négociés et celui où le dernier composant ouvre, l'exposition du porteur porte sur les
enchères d'ouverture auxquelles il peut participer, mais qu'il ne peut plus couvrir avec le
contrat arrivant à échéance.}\label{fig:m1:basis-roll-and-delivery:expiry}
\end{omfigure}

Un arbitragiste sur indice, acheteur du panier et vendeur du contrat à terme arrivant à
échéance, sort en vendant chaque action dans son enchère d'ouverture le matin de l'échéance~: le
produit de la vente est, par construction, le \omterm{def:m1:basis-roll-and-delivery:soq}{cours spécial d'ouverture}, et la
position disparaît sans risque de \omterm{def:m1:basis-roll-and-delivery:basis}{base}. C'est pourquoi le troisième vendredi
voit ses enchères d'ouverture compter parmi les plus importantes du trimestre
(\cref{ch:m1:auctions}).

\begin{example}[Moins trente-sept dollars]\label{ex:m1:basis-roll-and-delivery:wti}
Le contrat à terme WTI de mai 2020 devait arriver à échéance le 21 avril. Le 20 avril, il a ouvert
la séance à 17.73~\$~; entre environ 14 h 08 et la fin de la période de règlement-livraison
à 14 h 30 heure de New York, il s'est échangé sous zéro, pour la première fois depuis
son admission à la cote en 1983, a touché $-\$40.32$ et a réglé à
$-\$37.63$. Toutes les autres échéances ont réglé à un prix positif. Le rapport intermédiaire du régulateur
invoque un marché en surproduction, l'effondrement de la demande lié à la
pandémie, l'inquiétude sur le stockage au point de livraison, et relève que
environ un cinquième du volume du jour sur ce contrat a été échangé \emph{at
settlement}, à un prix connu seulement après 14 h 30. Un acheteur qui ne peut pas prendre
livraison doit vendre avant l'échéance, à n'importe quel prix qui se forme sur le marché~; la convergence
du contrat à terme et du comptant est assurée par ceux qui peuvent livrer et stocker, et ce
jour-là, peu d'entre eux le pouvaient.
\end{example}

\section{Négocier à un prix encore inconnu}

\begin{definition}[Exchange for physical]\label{def:m1:basis-roll-and-delivery:efp}
Un \emph{exchange for physical}\index{exchange for physical} (EFP) est une
transaction négociée de gré à gré, déclarée à la bourse, dans laquelle une partie
achète le sous-jacent au comptant et vend le contrat à terme, l'autre partie faisant
l'inverse, à un écart de prix convenu. Elle déplace une position entre le marché au comptant
et le marché à terme, sans passer par aucun des deux carnets d'ordres.
\end{definition}

\begin{definition}[Trade at settlement et basis trade at index close]\label{def:m1:basis-roll-and-delivery:tas}
Un \emph{trade at settlement}\index{trade at settlement} (TAS) est une transaction sur contrat à terme
exécutée pendant la séance à un prix égal au cours de règlement-livraison du jour,
augmenté ou diminué d'un nombre convenu de ticks, lequel n'est connu qu'après le
règlement-livraison.
Un \emph{basis trade at index close}\index{basis trade at index close} (BTIC) est une transaction sur contrat à terme à un prix égal au niveau de clôture
officiel de l'indice du jour, augmenté d'une \omterm{def:m1:basis-roll-and-delivery:basis}{base} convenue.
\end{definition}

Les deux répondent au même besoin~: un fonds valorisé à la clôture veut des contrats à terme à la
clôture, ou à une distance fixe de celle-ci, sans négocier dans les trente dernières
secondes. Une cotation BTIC est une cotation directe de la \omterm{def:m1:basis-roll-and-delivery:basis}{base}, et donc du \omterm{ex:m1:basis-roll-and-delivery:43}{taux de financement implicite}~:
c'est là que le \omterm{def:m1:delta-one-instruments:swap}{spread de financement} des contrats à terme sur indice s'observe le plus
nettement. Les deux concentrent aussi le risque dans la fenêtre de règlement-livraison~: les
contreparties des ordres TAS se couvrent sur le marché avant le règlement-livraison, et
le 20 avril 2020 cette couverture a pesé sur un marché sans acheteurs.

\section{Tutoriel~: juste valeur avec dividendes discrets}\label{tut:m1:basis-roll-and-delivery:fv}

\begin{tutorial}
\textbf{Objectif.} Calculer la juste valeur d'un contrat à terme sur indice à partir de dividendes
datés, l'inverser pour obtenir le \omterm{ex:m1:basis-roll-and-delivery:43}{taux de financement implicite}, tracer la \omterm{def:m1:basis-roll-and-delivery:basis}{base} sur la durée de
vie du contrat et la bande d'arbitrage. \textbf{État final~:} les trois
figures de données de ce chapitre.
\begin{tutsteps}
\item \textbf{Juste valeur}, chaque dividende étant réinvesti jusqu'à l'échéance.
\omcode{code/firm/fairvalue/firm_fairvalue.py}{21}{27}{Dividendes reportés jusqu'à l'échéance, et juste valeur.}\label{lst:m1:basis-roll-and-delivery:fv}
\item \textbf{Taux implicite} par dichotomie, car le réinvestissement des dividendes
dépend aussi du taux.
\omcode{code/firm/fairvalue/firm_fairvalue.py}{30}{42}{Le taux de financement auquel le prix de marché est juste.}\label{lst:m1:basis-roll-and-delivery:implied}
\item \textbf{La bande}, asymétrique lorsque le panier est cher à emprunter.
\omcode{code/firm/fairvalue/firm_fairvalue.py}{53}{62}{Cours des contrats à terme au-delà desquels l'arbitrage sur indice est rentable.}\label{lst:m1:basis-roll-and-delivery:band}
\item \textbf{Une position roulée} sur une courbe de matières premières.
\omcode{code/markets-1/21-basis-roll-and-delivery/python/basis_demo.py}{46}{53}{Valeur d'une position longue roulée à chaque période.}\label{lst:m1:basis-roll-and-delivery:rolled}
\end{tutsteps}
\textbf{Pour aller plus loin.} Rendre les dividendes incertains~: décaler chacun de
$\pm 10\%$ et voir quelle part de la bande cela consomme pour un contrat à terme à un an.
C'est pourquoi les contrats à terme sur indice à longue échéance et les contrats à terme sur dividendes existent comme
marchés distincts (\cref{ch:m1:index-and-single-stock-futures-worldwide}).
\end{tutorial}

\section{Construction~: le calculateur de juste valeur}\label{bld:m1:basis-roll-and-delivery:fv}

\begin{build}
\bfield{Rôle} Le \omterm{def:m1:what-a-trading-firm-does:market-maker}{teneur de marché} sur contrats à terme de la petite firme cote autour
de la juste valeur, sa stratégie d'arbitrage sur indice opère contre elle, et son système de risque
l'utilise pour décomposer le P\&L des contrats à terme en indice, carry et mauvaise valorisation.
\bfield{Interface} \texttt{Dividend(pay\_date, points)};
\texttt{fair\_value(spot, rate, today, expiry, divs)};
\texttt{implied\_rate(future, spot, today, expiry, divs)};
\texttt{ArbCosts}; \texttt{arbitrage\_band(\ldots)} renvoyant les prix sans arbitrage inférieur et
supérieur~; \texttt{roll\_richness\_bp(spread\_market, spot,
rate, today, near, far, divs)}. \bfield{Règles} Intérêt simple, actual/360~; dividendes comptés s'ils sont payés après
aujourd'hui et au plus tard à l'échéance~; dates d'échéance issues du référentiel des contrats
(\cref{bld:m1:futures-contracts-and-exchanges:master})~; points de dividende issus
de l'ajusteur d'opérations sur titres et des pondérations de l'indice
(\cref{bld:m1:shares-corporate-actions-indices:adjuster})~; \omterm{def:m1:financing:seclending}{frais d'emprunt} issus du
livre des emprunts (\cref{bld:m1:stock-loan-and-short-selling:borrow}).
\bfield{Tests d'acceptation} \texttt{code/firm/fairvalue/tests/}~: intérêt simple
sans dividende~; l'exemple du chapitre calculé à la main~; le taux implicite
comme inverse~; la bande asymétrique~; un \omterm{def:m1:basis-roll-and-delivery:roll}{roll} de 30 points de base trop cher.
\bfield{Pour aller plus loin} Une structure par termes de taux au lieu d'un $r$~; dividendes dans
la devise de l'indice pour un contrat à terme quanto~; une convention en date ex, et non en date de
paiement, avec l'écart signalé.
\end{build}

\begin{omsources}
\item CME, Rulebook Chapter 358, règles 35802.G et 35803.A (fin de la
négociation et règlement-livraison final), tel que déposé auprès de la CFTC.
\item US Commodity Futures Trading Commission, \emph{Interim Staff Report:
Trading in NYMEX WTI Crude Oil Futures Contract Leading up to, on, and around
April 20, 2020}, novembre 2020.
\item US Commodity Futures Trading Commission, \emph{Futures Glossary}
(exchange for physicals); CME Group, \emph{Basis Trade at Index Close (BTIC)}.
\item Eurex, \emph{EURO STOXX 50 Index Futures}, page produit.
\end{omsources}

\section{Exercices}

\begin{exercise}[$\star$]\label{exo:m1:basis-roll-and-delivery:1}
Avec $S = 6\,000$, $r = 4.2\%$, 91 jours et sans dividende, donner la juste valeur
et la \omterm{def:m1:basis-roll-and-delivery:basis}{base}.
\end{exercise}

\begin{exercise}[$\star$]\label{exo:m1:basis-roll-and-delivery:2}
Ajouter les trois dividendes de l'\cref{ex:m1:basis-roll-and-delivery:43}, versés 64,
32 et 8 jours avant l'échéance. Donner leur valeur à l'échéance et la \omterm{def:m1:basis-roll-and-delivery:carry}{fair value}.
\end{exercise}

\begin{exercise}[$\star$]\label{exo:m1:basis-roll-and-delivery:3}
Une courbe de pétrole brut a ses deux premiers contrats mensuels à 70.00 et 70.84.
Nommer la forme de la courbe et donner le rendement de \omterm{def:m1:basis-roll-and-delivery:roll}{roll} annualisé d'une position longue. Idem
pour 70.00 et 69.30.
\end{exercise}

\begin{exercise}[$\star\star$]\label{exo:m1:basis-roll-and-delivery:4}
Avec les coûts de l'\cref{ex:m1:basis-roll-and-delivery:band}, le contrat à terme se
traite à 6\,050.50. Décrire l'arbitrage, jambe par jambe, et donner le
profit verrouillé par contrat. Qu'est-ce qui pourrait encore mal tourner~?
\end{exercise}

\begin{exercise}[$\star\star$]\label{exo:m1:basis-roll-and-delivery:5}
Le contrat à terme de l'\cref{ex:m1:basis-roll-and-delivery:43} se traite à 6\,047.20.
Vérifier le \omterm{ex:m1:basis-roll-and-delivery:43}{taux de financement implicite}. Si le taux de référence monte de 25 points de base
demain, l'indice restant inchangé, de combien la \omterm{def:m1:basis-roll-and-delivery:carry}{fair value} varie-t-elle~?
\end{exercise}

\begin{exercise}[$\star\star$]\label{exo:m1:basis-roll-and-delivery:6}
Un vendredi d'échéance, l'indice publié ouvre à 5\,990, se traite entre
5\,985 et 6\,020 dans la journée, et le \omterm{def:m1:basis-roll-and-delivery:soq}{cours spécial d'ouverture}
vaut 5\,982. Expliquer comment ce cours peut se situer sous le plus bas de la journée. Qui s'en soucie~?
\end{exercise}

\begin{exercise}[$\star\star\star$]\label{exo:m1:basis-roll-and-delivery:7}
\emph{Codage.} Avec \texttt{rolled\_index}, reproduire les deux valeurs finales de la
\cref{fig:m1:basis-roll-and-delivery:rolled}. Puis faire décroître linéairement le \omterm{def:m1:basis-roll-and-delivery:contango}{contango}
de 1.2~\% à zéro sur les 24 mois et donner la valeur finale.
\end{exercise}

\begin{exercise}[$\star\star\star$]\label{exo:m1:basis-roll-and-delivery:8}
\emph{Trouver la faille.} «~Les contrats à terme sur le S\&P se traitent 43 points au-dessus de l'indice
ce matin~: les traders sur contrats à terme anticipent une hausse du marché.~» Corriger
l'affirmation. Que \emph{révéleraient} les contrats à terme avant l'ouverture du marché
boursier~?
\end{exercise}

\section{Problème~: Faire rouler un milliard}

\begin{problem}[{Problème du week-end --- le coût trimestriel de rester long}]\label{pb:m1:basis-roll-and-delivery:1}
Un fonds détient une exposition de 1~milliard de dollars au S\&P 500 via des E-mini (indice 6\,000,
multiplicateur 50~\$), sa trésorerie étant placée au taux de référence de 4.2~\%. Nous sommes le 18
septembre 2026~; il détient le contrat de décembre (échéance le 18 décembre) et passera à
mars (19 mars 2027). Les dividendes entre les deux échéances sont de 6.0,
9.5 et 5.5 points, versés les 15 janvier, 15 février et 10 mars.

\medskip\noindent\textbf{Partie I --- Le \omterm{def:m1:basis-roll-and-delivery:roll}{roll} équitable.}
\begin{enumerate}
\item Combien de contrats le fonds détient-il~?
\item Donner la \omterm{def:m1:basis-roll-and-delivery:carry}{fair value} du contrat de décembre.
\item Donner la \omterm{def:m1:basis-roll-and-delivery:carry}{fair value} du contrat de mars.
\item Donner le \omterm{def:m1:matching-algorithms-and-implied-spreads:calendar}{calendar spread} équitable (mars moins décembre).
\item Le décomposer en intérêts et dividendes.
\end{enumerate}

\medskip\noindent\textbf{Partie II --- Le \omterm{def:m1:basis-roll-and-delivery:roll}{roll} de marché.}
Le spread se traite 4.55 points au-dessus de la \omterm{def:m1:basis-roll-and-delivery:carry}{fair value}.
\begin{enumerate}[resume]
\item Donner la cherté en points de base par an.
\item Donner son coût en dollars pour ce \omterm{def:m1:basis-roll-and-delivery:roll}{roll}.
\item Le fonds paie un demi-tick de spread (le tick vaut 0.05) à l'exécution.
Donner ce coût.
\item Donner le coût annuel total du \omterm{def:m1:basis-roll-and-delivery:roll}{roll}, en dollars et en points de base
de l'exposition.
\item Comparer avec les 7 points de base par an d'un fonds indiciel
(\cref{ch:m1:delta-one-instruments}). Qu'est-ce qui doit encore entrer dans la comparaison~?
\end{enumerate}

\medskip\noindent\textbf{Partie III --- Quand rouler.}
\begin{enumerate}[resume]
\item La plus grande partie des \omterm{def:m1:futures-contracts-and-exchanges:oi}{positions ouvertes} roule la semaine précédant l'échéance. Quel est
l'argument pour rouler avec la foule, et pour rouler plus tôt~?
\item En décembre, la cherté est généralement la plus élevée. Pourquoi~?
\item Le fonds pourrait vendre la \omterm{def:m1:basis-roll-and-delivery:basis}{base} via le BTIC à la place. Qu'apprend-il d'une cotation BTIC
qu'il n'apprend pas du cours du contrat à terme~?
\item Une banque propose de prendre la contrepartie du \omterm{def:m1:basis-roll-and-delivery:roll}{roll} à la \omterm{def:m1:basis-roll-and-delivery:carry}{fair value} plus 20
points de base. Pourquoi le peut-elle, et pourquoi le ferait-elle~?
\end{enumerate}

\medskip\noindent\textbf{Partie IV --- Jugement.}
\begin{enumerate}[resume]
\item Un fonds de matières premières fait face à la même décision sur une courbe en
\omterm{def:m1:basis-roll-and-delivery:contango}{contango} prononcé. Qu'est-ce qui diffère~?
\item Un fonds qui oublie de rouler une position E-mini la conserve jusqu'à l'échéance.
Que se passe-t-il~? Et si le contrat était un WTI~?
\item Pourquoi la cherté du \omterm{def:m1:basis-roll-and-delivery:roll}{roll} est-elle la mesure de quelque chose de plus large que
les coûts de ce fonds~?
\item Qui gagne les 30 points de base~?
\item Énoncer le \emph{résultat nommé}~: le coût du \omterm{def:m1:basis-roll-and-delivery:roll}{roll} en points de base par
an.
\item En une phrase~: pourquoi le contrat à terme converge-t-il vers l'indice~?
\end{enumerate}
\end{problem}

\section{Questions d'entretien}

\begin{interviewq}[$\star$ \iqroles{trader, chercheur, banque}]\label{iq:m1:basis-roll-and-delivery:1}
Pourquoi un contrat à terme sur indice boursier se traite-t-il avec une prime sur l'indice~? Quand
se traiterait-il avec une décote~?
\end{interviewq}

\begin{interviewq}[$\star$ \iqroles{trader, chercheur}]\label{iq:m1:basis-roll-and-delivery:2}
Définissez le \omterm{def:m1:basis-roll-and-delivery:contango}{contango} et la \omterm{def:m1:basis-roll-and-delivery:contango}{backwardation}. Le \omterm{def:m1:basis-roll-and-delivery:contango}{contango} signifie-t-il que le marché anticipe
une hausse des prix~?
\end{interviewq}

\begin{interviewq}[$\star\star$ \iqroles{trader, chercheur}]\label{iq:m1:basis-roll-and-delivery:3}
Décrivez-moi un arbitrage sur indice, y compris la manière dont vous en sortez.
\end{interviewq}

\begin{interviewq}[$\star\star$ \iqroles{chercheur, ingénieur ML}]\label{iq:m1:basis-roll-and-delivery:4}
Vous construisez une série continue de prix de contrats à terme pour la recherche. Quelles sont
vos options à chaque \omterm{def:m1:basis-roll-and-delivery:roll}{roll}, et à quoi sert chacune~?
\end{interviewq}

\begin{interviewq}[$\star\star$ \iqroles{trader, banque}]\label{iq:m1:basis-roll-and-delivery:5}
Comment un contrat à terme sur pétrole brut a-t-il pu se régler à un prix négatif, et cela pourrait-il arriver
à un contrat à terme sur indice boursier~?
\end{interviewq}

\begin{interviewq}[$\star\star\star$ \iqroles{trader, chercheur}]\label{iq:m1:basis-roll-and-delivery:6}
Le \omterm{def:m1:basis-roll-and-delivery:roll}{roll} du S\&P se traite 60 points de base cher sur décembre. Comment le traitez-vous, et
qu'est-ce qui vous empêche de le faire en taille illimitée~?
\end{interviewq}
